% J1.6: recorrido de los datos tabulares, de los ficheros a las figuras y de vuelta a un fichero.
\documentclass[tikz,border=4pt]{standalone}
\usepackage{resonancia}
\begin{document}
\begin{tikzpicture}
  \node[cajaNaranja, minimum width=2.4cm] (csv) at (0,1.0) {\textbf{CSV}\\[-1pt]{\footnotesize \texttt{CSV.read}}};
  \node[cajaNaranja, minimum width=2.4cm] (xl) at (0,-0.6) {\textbf{Excel}\\[-1pt]{\footnotesize \texttt{XLSX.readtable}}};
  \node[cajaAzul, minimum width=3.0cm, minimum height=1.6cm] (df) at (4.2,0.2) {\textbf{DataFrame}\\[-1pt]{\footnotesize columnas con tipo}\\[-1pt]{\footnotesize \texttt{missing}}};
  \node[cajaVerde, minimum width=2.4cm] (mat) at (8.4,1.0) {\textbf{Matriz} $N\times c$\\[-1pt]{\footnotesize \texttt{Matrix(df[:, 2:4])}}};
  \node[cajaVerde, minimum width=2.4cm] (fig) at (8.4,-0.6) {\textbf{Figura}\\[-1pt]{\footnotesize CairoMakie}};
  \node[caja, draw=rVioleta, fill=rVioleta!8, minimum width=2.4cm] (out) at (4.2,-2.3) {\textbf{CSV / Excel}\\[-1pt]{\footnotesize \texttt{CSV.write}, \texttt{XLSX.writetable}}};
  \draw[flecha] (csv) -- (df); \draw[flecha] (xl) -- (df);
  \draw[flecha] (df) -- (mat); \draw[flecha] (df) -- (fig);
  \draw[flecha, color=rVioleta] (df) -- (out);
\end{tikzpicture}
\end{document}
